\documentclass[11pt,a4paper]{article}
\usepackage[T1]{fontenc}
\usepackage{lmodern,amsmath,amssymb,booktabs,microtype}
\usepackage[margin=24mm]{geometry}
\usepackage[colorlinks=true,urlcolor=blue,linkcolor=blue]{hyperref}
\newcommand{\Kop}{\mathsf K}
\newcommand{\Aop}{\mathcal A}
\DeclareMathOperator{\sech}{sech}
\title{A concrete bulk mass target\\Connecting the singlet programme to the 't Hooft equation}
\author{Alan Day \and Chengheng\thanks{AI research assistant (OpenAI).}}
\date{Local research note --- 6 November 2026}
\begin{document}\maketitle
\begin{abstract}
We connect the current singlet programme to the precise planar meson
problem selected by the proposed Susskind parameter dictionary. The
target has zero renormalized quark mass, not zero bare mass. Its known
mass gap and first six levels are reproduced in two independent integral
representations. A quadratic-form comparison controls the approach to
this target at finite interaction order. We then formulate a spatial
response test that would distinguish a single meson from a mixture,
while retaining the previously computed boost correlators.
\end{abstract}
\section{The mass parameter selected by the proposal}
Write $\tau=g_{\rm YM}^2N$ for the bulk tension scale and distinguish
the renormalized quark mass $m_R$ from the bare mass $m_b$.
The proposed dictionary in \cite{MSS26,MSS25} gives
\begin{equation}
 m_R^2=m_b^2-\frac{\tau}{\pi},\qquad
 \frac{\tau}{m_R^2}=p^2,\qquad
 \lambda_{\rm SYK}=\frac{p^2}{N}.
\end{equation}
In the fixed-tension limit, $m_R^2=\tau/p^2\to0$ but
$m_b^2\to\tau/\pi$. This is the special case solved in detail by
Fateev--Lukyanov--Zamolodchikov (FLZ) \cite{FLZ}.
To avoid the two different uses of lambda, denote their spectral
parameter by $\nu$ and their dimensionless quark parameter by $a$:
\begin{equation}
 a=\frac{\pi m_b^2}{\tau}-1=\frac{\pi}{p^2},\qquad
 M_r^2=2\pi\tau\nu_r,\qquad r=0,1,\ldots.
\end{equation}
The limiting target is $a=0$. The bare-massless, chiral problem is
$a=-1$ and has different infrared physics.

This is a bulk planar spectral target. Using it as the leading small
$\lambda_{\rm SYK}$ spectrum remains a matching task: it is not a
derivation of a microscopic spatial generator or of finite-coupling
meson interactions. The original proposal itself identifies quantitative
spectral agreement as a test still to be performed.

\newpage
\section{The actual internal mass operator}
On $L^2((0,1),dx)$ use FLZ's finite-part equation
\begin{equation}
 \Kop_a\varphi(x)=\frac{a}{x(1-x)}\varphi(x)
 -\operatorname{FP}\!\int_0^1\frac{\varphi(y)}{(y-x)^2}\,dy
 =2\pi^2\nu\varphi(x).\label{eq:thooft}
\end{equation}
Here $x$ is a quark momentum fraction \emph{inside one meson}.
It is not the fraction shared by the two composite singlets in the
preceding ordering-metric note. All masses in this section are bulk
Minkowski masses; SYK time instead generates a boost.

Subtracting the diagonal singularity yields the positive form at $a=0$,
\begin{align}
 q_0[f]&=\frac12\int_0^1\!\!\int_0^1
       \frac{|f(x)-f(y)|^2}{(x-y)^2}\,dx\,dy
       +\int_0^1\frac{|f(x)|^2}{x(1-x)}\,dx,\label{eq:form}\\
 q_a[f]&=q_0[f]+a\int_0^1\frac{|f(x)|^2}{x(1-x)}\,dx.
\end{align}
The positive closure defines the self-adjoint operator. The endpoint
potential alone gives $q_0[f]\ge4\|f\|^2$. Thus setting $m_R=0$
does not remove the mass gap. At $a=-1$ that endpoint term cancels;
the constant solution has zero mass in the chiral form domain.

\paragraph{A controlled approach to the bulk target.}
Since the last integral in \eqref{eq:form} is at most $q_0[f]$,
the closed forms have the same domain for $a\ge0$ and obey
\begin{equation}
 q_0\le q_a\le(1+a)q_0.
\end{equation}
The min--max principle, applied to the known discrete spectrum, gives
\begin{equation}
 \boxed{0\le\frac{\nu_r(a)}{\nu_r(0)}-1\le a=\frac{\pi}{p^2}.}
 \label{eq:bound}
\end{equation}
This is a form-comparison consequence of the standard equation, not a
new meson solution. It controls every fixed level, indeed the relative
bound is independent of $r$, in this planar family at fixed $\tau$.
It controls neither the SYK double-scaling error nor nonplanar corrections.

\paragraph{Where the previous translation test applies.}
The earlier $\Aop=C_{\rm tot}-C_1-C_2$ acts on two \emph{composite}
singlets. Replacing its relative variable by the internal quark variable
in \eqref{eq:thooft} has no justification. Likewise the insertion vertex
$i\sqrt{h}\sqrt{\Aop}/2$ has not been shown to be an on-shell transition
map between invariant sectors. Its conditional free-translation
incompatibility remains valid, but does not test the operator
\eqref{eq:thooft} by itself.

\newpage
\section{Reproducing the known spectrum}
We use the Chebyshev method cited in FLZ Section 5. Put $2x-1=\cos\theta$
and expand $\varphi=\sum_{n\ge1}c_n\sin(n\theta)$. The finite-part identity is
\begin{equation}
 -\operatorname{FP}\!\int_0^1
 \frac{\sin(n\arccos(2y-1))}{(y-x)^2}\,dy
 =2\pi n\,U_{n-1}(2x-1).
\end{equation}
Let $D_{mn}=n\delta_{mn}$. Projecting the equation gives
\begin{equation}
 Dc=2\nu Bc,\qquad
 B_{mn}=\begin{cases}
 \displaystyle\frac1{1-(m-n)^2}-\frac1{1-(m+n)^2},&m+n\text{ even},\\
 0,&m+n\text{ odd}.
 \end{cases}
\end{equation}
The matrix $2D^{-1/2}BD^{-1/2}$ has eigenvalues $1/\nu$.
The physical norm is $c^*Bc/2$, not the unweighted coefficient norm.
Reflection $x\mapsto1-x$ acts on mode $n$ by $(-1)^{n+1}$.

\begin{center}
\begin{tabular}{rrrr}
\toprule
$r$ & $2\nu_r$ (reproduced) & $M_r^2/\tau=2\pi\nu_r$ & reflection\\
\midrule
0 & 0.737061746293 & 2.315548 & $+$\\
1 & 1.753731336918 & 5.509509 & $-$\\
2 & 2.748160912371 & 8.633602 & $+$\\
3 & 3.751057581705 & 11.784295 & $-$\\
4 & 4.749295381038 & 14.920351 & $+$\\
5 & 5.750492623649 & 18.065705 & $-$\\
\bottomrule
\end{tabular}
\end{center}
The first numerical column agrees with FLZ Tables 1 and 2; it is not an
independent prediction of DSSYK. The labels are FLZ labels, starting
at zero. Identifying individual derivative operators with these
eigenstates requires their overlaps, not just matching their ordering.

For an analytic check, $B\le(\pi/2)I$ gives
$\Kop_0\ge2\pi$ in the sine basis. The trial $\sin\theta$ gives
$\nu_0\le3/8$. Hence
\begin{equation}
 2\tau\le M_0^2\le\frac{3\pi}{4}\tau,
\end{equation}
consistent with the reproduced gap. These are elementary variational
bounds, not numerical error bars on the tabulated eigenvalues.

\newpage
\section{A spatial observable that can test the dictionary}
Given the bulk spectrum, the one-meson space can be represented as
$L^2(\mathbb R_+,dp)\otimes\ell^2(r)$. With dimensionless $p$ and an
energy unit $\mu>0$, define
\begin{equation}
 P^+=\mu p,\qquad P^-|r\rangle=\frac{M_r^2}{\mu p}|r\rangle,
 \qquad B=i(p\partial_p+\tfrac12).
\end{equation}
These standard positive-energy generators obey the Poincare relations.
They are now supplied with masses from a specific bulk equation, rather
than an arbitrary function of a boost representation label.

Take the explicitly separable candidate wave packet
\begin{equation}
 |a\rangle=\psi_d(p)\sum_r c_{ar}|r\rangle,\qquad
 \psi_d(p)=\frac{2^d}{\sqrt{\Gamma(2d)}}p^{d-1/2}e^{-p},\qquad
 \sum_r|c_{ar}|^2=1.
\end{equation}
Every such packet has
$\langle a|e^{-2iuB}|a\rangle=\sech^{2d}u$, regardless of its mass
weights. Thus the known singlet time kernels alone cannot select the
coefficients $c_{ar}$.

The opposite-null translation probes the mass distribution; $s$ has
units of inverse energy:
\begin{equation}
 F_a(s)=\langle a|e^{-isP^-}|a\rangle
 =\sum_r|c_{ar}|^2\int_0^\infty |\psi_d(p)|^2
                 e^{-isM_r^2/(\mu p)}\,dp.\label{eq:response}
\end{equation}
For $d>1$ and finite fourth mass moment,
\begin{align}
 \langle M^2\rangle_a&=\frac{i\mu(2d-1)}2F_a'(0),\\
 \langle M^4\rangle_a&=-\frac{\mu^2(2d-1)(2d-2)}4F_a''(0).
\end{align}
In particular
\begin{equation}
 \boxed{\mathcal R_a=\frac{2d-2}{2d-1}
              \frac{-F_a''(0)}{|F_a'(0)|^2}
       =\frac{\langle M^4\rangle_a}{\langle M^2\rangle_a^2}\ge1.}
\end{equation}
Equality means support at one mass, within this separable construction.
This supplies a concrete test of the proposed lowest-singlet/lowest-meson
identification. A microscopic definition of the translation in
\eqref{eq:response} is still required; inserting bulk values by hand
would only define a model, not establish a duality.

\newpage
\appendix
\section{Independent checks and the next construction}
FLZ Appendix B gives another representation, with
\begin{equation}
 K(t,t')=\frac{(t-t')/\sinh(t-t')}
 {\sqrt{t\coth t}\sqrt{t'\coth t'}},\qquad
 K\chi=\nu^{-1}\chi.
\end{equation}
We discretize it with weighted Gauss--Legendre quadrature on $[-L,L]$.
The six-level maximum relative errors against the published values are
\begin{center}
\begin{tabular}{rrr}
\toprule
$L$ & nodes & maximum relative error\\\midrule
16&160&$1.5020\,10^{-2}$\\
24&240&$2.9125\,10^{-5}$\\
24&360&$2.9125\,10^{-5}$\\
32&480&$8.1708\,10^{-9}$\\
40&600&$8.0381\,10^{-13}$\\
40&800&$7.8071\,10^{-13}$\\\bottomrule
\end{tabular}
\end{center}
Increasing nodes at fixed $L=24$ does not remove the tail error. The
initial $L=32$ run failed the planned $2\,10^{-9}$ threshold; the record
was retained and the domain enlarged without relaxing the threshold.

Chebyshev sizes $32,64,128,256$ give maximum reference errors
$1.51\,10^{-8}$, $2.81\,10^{-10}$, $4.74\,10^{-12}$ and
$8.80\,10^{-14}$. The final two representations differ relatively by
$6.93\,10^{-13}$. Saved physical Gram and parity checks use the
correct $dx$ measure. At 65 digits, twelve subtracted finite-part
integrals, 36 angular Gram entries and four inverse-momentum moments
check the analytic formulas. The initial hypergeometric Chebyshev
evaluation failed; a finite polynomial recurrence fixed that evaluation.
Eight finite-basis comparisons at $p=4,8,16,32$ also check the form
bound. All original failures are retained. These controls are not certified
continuum eigenvalue enclosures.

The next task is to identify a microscopic spatial insertion or a
bilocal-to-meson reconstruction kernel. It must reproduce the singlet
boost data and the sourced mass operator together, and then the
nonzero multiplication data. The one-meson internal coordinate, the
two-meson relative coordinate, and the boost wave-packet label must
remain distinct in that construction.

\begin{thebibliography}{9}
\bibitem{MSS26} S. Miyashita, Y. Sekino and L. Susskind,
\emph{Holograms and Standard Models},
\href{https://arxiv.org/abs/2607.05678v2}{arXiv:2607.05678v2},
Secs.~4.1--4.3.
\bibitem{MSS25} S. Miyashita, Y. Sekino and L. Susskind,
\emph{DSSYK at Infinite Temperature: The Flat-Space Limit and the
't Hooft Model}, \href{https://arxiv.org/abs/2506.18054v3}{arXiv:2506.18054v3},
Secs.~5.2--5.3 and 6.2.
\bibitem{FLZ} V. A. Fateev, S. L. Lukyanov and A. B. Zamolodchikov,
\emph{On mass spectrum in 't Hooft's 2D model of mesons},
\href{https://arxiv.org/abs/0905.2280v2}{arXiv:0905.2280v2},
Eqs.~(1.1)--(1.5), Sec.~5 and Appendix B.
\end{thebibliography}
\end{document}
